\chapter{Теория}

\begin{figure}[H]
    \centering
    \includegraphics[width=0.75\textwidth]{protos.png}
    \caption{Протоколы OSI и TCP/IP}
\end{figure}

\section{Перечень протоколов}

\subsection{Доступ к сети (L1 и L2 в модели OSI)}

\subsubsection{PPP (Point-to-Point Protocol)}
PPP — двухточечный протокол канального уровня сетевой модели OSI (уровень доступа к
сети в стеке TCP/IP). Обычно используется для установления прямой связи между двумя
узлами сети, причём он может обеспечить аутентификацию соединения, шифрование и сжатие
данных. RFC: 1661.

\subsubsection{PPPoE (Point-to-point protocol over Ethernet)}
PPPoE — сетевой протокол канального уровня модели OSI (уровень доступа к сети в стеке
TCP/IP) передачи кадров PPP через Ethernet. В основном используется xDSL-сервисами.
Предоставляет дополнительные возможности (аутентификация, сжатие данных, шифрование).
RFC: 2516.

\subsubsection{ARP (Address Resolution Protocol)}
ARP — протокол (уровень доступа к сети в стеке TCP/IP) в компьютерных сетях,
предназначенный для определения MAC-адреса другого компьютера по известному IP-адресу.
RFC: 826.

\subsubsection{DSL (Digital Subscriber Line)}
Под DSL понимают технологии, повышающие пропускную способность обычного медного
телефонного кабеля, соединяющего клиента и провайдера Internet (ISP). Технология DSL
позволяет одновременно и на высокой скорости передавать голосовую информацию и данные
по обычной паре телефонных проводов. RFC: 4319, 3728.

\subsubsection{MPLS (Multiprotocol Label Switching)}
MPLS — механизм в высокопроизводительной телекоммуникационной сети,
осуществляющий передачу данных от одного узла сети к другому с помощью меток. MPLS
является масштабируемым и независимым от каких-либо протоколов механизмом передачи
данных. RFC: 3031.

\subsubsection{ATM (Asynchronous Transfer Mode)}
ATM — сетевая высокопроизводительная технология коммутации и мультиплексирования
пакетов. Пакеты представляют собой ячейки фиксированного размера в 53 байта, где первые
5 байт используются под заголовок. Является разновидностью быстрой коммутации пакетов.
RFC: 1932.

\subsubsection{Ethernet IEEE 802.2}
Ethernet IEEE 802.2 является исходным названием стандарта ISO / IEC 8802-2, который
определяет управление логическим каналом (LLC) как верхнюю часть уровня канала передачи
данных модели OSI (уровень доступа к сети в стеке TCP/IP). Первоначальный стандарт,
разработанный Институтом инженеров по электротехнике и электронике (IEEE) в
сотрудничестве с Американским национальным институтом стандартов (ANSI), был принят
Международной организацией по стандартизации (ISO) в 1998 году, но он все ещё остается
неотъемлемой частью семейства из IEEE 802 стандартов для локальных и городских сетей. LLC
— это программный компонент, который обеспечивает единый интерфейс для пользователя
службы передачи данных, обычно сетевого уровня. RFC: 1700.

\subsubsection{Ethernet IEEE 802.3}
Ethernet IEEE 802.3 — семейство технологий пакетной передачи данных между
устройствами для компьютерных и промышленных сетей. Стандарты Ethernet определяют
проводные соединения и электрические сигналы на физическом уровне, формат кадров и
протоколы управления доступом к среде на канальном уровне модели OSI (уровень доступа к
сети в стеке TCP/IP). Ethernet, в основном, описывается стандартами IEEE группы 802.3.
Ethernet стал одной из самых распространённых технологий ЛВС в середине 1990-х годов,
вытеснив такие устаревшие технологии как Token Ring, FDDI и ARCNET.

\subsubsection{Token Bus IEEE 802.4}
Стандарт 802.4 описывает широкополосную шину локальной сети на основе коаксиального
кабеля, в которой используется маркер для управления доступом к носителю. 802.4 является
частью набора протоколов автоматизации производства (Manufacturing Automation Protocol),
предлагаемых для использования в промышленных условиях. Сигналы в широкополосном
коаксиальном кабеле не подвержены влиянию электромагнитных помех, связанных с
промышленным производством. RFC: 1230. 

\subsubsection{Token Ring IEEE 802.5}
Token Ring — протокол передачи данных в локальной вычислительной сети (LAN) с
топологией кольца и «маркерным доступом». Находится на канальном уровне (DLL) модели
OSI (уровень доступа к сети в стеке TCP/IP). Станции в локальной вычислительной сети Token
Ring логически организованы в кольцевую топологию с данными, передаваемыми
последовательно от одной станции в кольце к другой. RFC: 1231.

\subsubsection{FDDI (Fiber Distributed Data Interface)}
Технология FDDI — это первая технология локальных сетей, в которой средой передачи
данных является волоконно-оптический кабель. Протокол рассчитан на физическую скорость
передачи информации 100 Мбит/с и предназначен для сетей с суммарной длиной до 100 км (40
км для мультимодовых волокон) при расстоянии между узлами 2 км или более. RFC: 1390.

\subsubsection{ISDN (Integrated Services Digital Network)}
ISDN — технология, обеспечивающая передачу цифрового сигнала по телефонным каналам
с предоставлением различных услуг. Основное назначение ISDN — передача данных со
скоростью до 64 кбит/с по абонентской проводной линии и обеспечение интегрированных
телекоммуникационных услуг (телефон, факс и пр.). RFC: 4715.

\subsubsection{Wi-Fi IEEE 802.11}
IEEE 802.11 — набор стандартов связи для коммуникации в беспроводной локальной
сетевой зоне частотных диапазонов 0,9; 2,4; 3,6; 5 и 60 ГГц. Пользователям более известен по
названию Wi-Fi, фактически являющемуся брендом, предложенным и продвигаемым
организацией Wi-Fi Alliance. Получил широкое распространение благодаря развитию
мобильных электронно-вычислительных устройств: КПК и ноутбуков.

\subsubsection{Bluetooth IEEE 802.15}
Bluetooth — технология беспроводной передачи данных, которая обеспечивает соединение
между собой устройств, имеющих встроенный Bluetooth-модуль.

\subsection{Интернет (L3 в модели OSI)}

\subsubsection{IPv4 (Internet Protocol version 4)}
IPv4 — четвёртая версия интернет-протокола (IP). Первая широко используемая версия.
IPv4 использует 32-битные (четырёхбайтные) адреса, ограничивающие адресное пространство
4{,}3 миллиардами возможными уникальными адресами.
Традиционной формой записи IPv4-адреса является запись в виде четырёх десятичных
чисел (от 0 до 255), разделённых точками. Через дробь указывается длина маски подсети.
RFC: 791.

\subsubsection{IPv6 (Internet Protocol version 6)}
IPv6 — новая версия интернет-протокола (IP), призванная решить проблемы, с которыми
столкнулась предыдущая версия (IPv4) при её использовании в Интернете, за счёт целого ряда
принципиальных изменений. Протокол был разработан IETF. Длина адреса IPv6 составляет 128
бит, в отличие от адреса IPv4, длина которого равна 32 битам. RFC: 8200.

\subsubsection{ICMP (Internet Control Message Protocol)}
ICMP — сетевой протокол, входящий в стек протоколов TCP/IP. В основном ICMP
используется для передачи сообщений об ошибках и других исключительных ситуациях,
возникших при передаче данных, например, запрашиваемая услуга недоступна, хост или
маршрутизатор не отвечают. RFC: 792.

\subsubsection{RIP (Routing Information Protocol)}
RIP — один из самых простых протоколов маршрутизации. Применяется в небольших
компьютерных сетях, позволяет маршрутизаторам динамически обновлять маршрутную
информацию, получая её от соседних маршрутизаторов. RFC: 2453.

\subsubsection{OSPF (Open Shortest Path First)}
OSPF — протокол динамической маршрутизации, основанный на технологии отслеживания
состояния канала, использующий алгоритм Дейкстры для нахождения кратчайшего пути.
Протокол OSPF был разработан IETF в 1988 году. RFC: 2328.

\subsubsection{NAT (Network Address Translation)}
NAT — это механизм в сетях TCP/IP, позволяющий преобразовывать IP-адреса транзитных
пакетов. Преобразование адреса методом NAT может производиться почти любым
маршрутизирующим устройством: маршрутизатором, сервером доступа, межсетевым экраном.
Наиболее популярным является SNAT, суть механизма которого состоит в замене адреса
источника при прохождении пакета в одну сторону и обратной замене адреса назначения в
ответном пакете. RFC: 3022.

\subsection{Транспортный (L4 в модели OSI)}

\subsubsection{TCP (Transmission Control Protocol)}
TCP — один из основных протоколов передачи данных интернета. Предназначен для
управления передачей данных Интернета. Пакеты в TCP называются сегментами.
Обеспечивает гарантированную передачу сегментов в интернете. RFC: 793.

\subsubsection{UDP (User Datagram Protocol)}
UDP — один из ключевых элементов набора сетевых протоколов для Интернета. С UDP
компьютерные приложения могут посылать сообщения другим хостам по IP-сети без
необходимости предварительного сообщения для установки специальных каналов передачи
данных. RFC: 768.

\subsubsection{RTCP (Real-Time Transport Control Protocol)}
RTCP — протокол, используемый совместно с RTP. RTCP базируется на периодической
передаче управляющих пакетов всем участникам сессии, используя тот же механизм рассылки,
что и для пакетов данных. Протокол RTCP используется для передачи информации о задержках
и потерях медиа-пакетов, джиттер-буфере, уровне звукового сигнала. RFC: 3550.

\subsubsection{SCTP (Stream Control Transmission Protocol)}
SCTP — протокол транспортного уровня в компьютерных сетях, появившийся в 2000 году в
IETF. Как и любой другой протокол передачи данных транспортного уровня, SCTP работает
аналогично TCP или UDP. Будучи более новым протоколом, SCTP имеет несколько
нововведений, таких как многопоточность, защита от DDoS-атак, синхронное соединение между
двумя хостами по двум и более независимым физическим каналам (multi-homing).
RFC: 4960, 3286.

\subsubsection{FCP (Fibre Channel Protocol)}
FCP — транспортный протокол (как TCP в IP-сетях), инкапсулирующий протокол SCSI по
сетям Fibre Channel. Является основой построения сетей хранения данных.

\subsubsection{DCCP (Datagram Congestion Control Protocol)}
DCCP — протокол транспортного уровня, разрабатываемый IETF. Принят в качестве
стандарта в марте 2006 года. Он предоставляет механизмы для отслеживания перегрузок в
сети, избегая необходимости создавать их на прикладном уровне. Этот протокол не гарантирует
доставку информации в нужном порядке. RFC: 4340.

\subsection{Уровень приложений (L5--L7 в модели OSI)}

\subsubsection{HTTP (HyperText Transfer Protocol)}
HTTP — протокол прикладного уровня передачи данных, изначально в виде гипертекстовых
документов в формате HTML, в настоящее время используется для передачи произвольных
данных. Обмен сообщениями идёт по обыкновенной схеме «запрос-ответ». Для идентификации
ресурсов HTTP использует глобальные URI. В отличие от многих других протоколов, HTTP не
сохраняет своего состояния. Это означает отсутствие сохранения промежуточного состояния
между парами «запрос-ответ». RFC: 1945, 2616, 7240.

\subsubsection{FTP (File Transfer Protocol)}
FTP — протокол передачи файлов по сети, появившийся в 1971 году задолго до HTTP и
даже до TCP/IP, благодаря чему является одним из старейших прикладных протоколов.
Изначально FTP работал поверх протокола NCP, на сегодняшний день широко используется
для распространения ПО и доступа к удалённым хостам. RFC: 959.

\subsubsection{POP3 (Post Office Protocol Version 3)}
POP3 — стандартный интернет-протокол прикладного уровня, используемый клиентами
электронной почты для получения писем с удалённого сервера по TCP-соединению.
RFC: 1939, 5034.

\subsubsection{SMTP (Simple Mail Transfer Protocol)}
SMTP — широко используемый сетевой протокол, предназначенный для передачи
электронной почты в сетях TCP/IP. RFC: 5321.

\subsubsection{RDP (Remote Desktop Protocol)}
RDP — проприетарный протокол прикладного уровня, использующийся для обеспечения
удалённой работы пользователя с сервером, на котором запущен сервис терминальных
подключений. RFC: 908, 1151.

\subsubsection{SMB (Server Message Block)}
SMB — сетевой протокол прикладного уровня (уровень приложений) для удалённого
доступа к файлам, принтерам и другим сетевым ресурсам, а также для межпроцессного
взаимодействия.

\subsubsection{RTP (Real-time Transport Protocol)}
RTP работает на уровне приложений и используется при передаче трафика реального
времени. Протокол RTP переносит в своём заголовке данные, необходимые для
восстановления аудиоданных или видеоизображений в приёмном узле, а также данные о типе
кодирования информации (JPEG, MPEG и т.\,п.). RFC: 4571.

\subsubsection{TELNET (Teletype Network)}
TELNET — сетевой протокол для реализации текстового терминального интерфейса по
сети (в современной форме — при помощи TCP). Протокол TELNET использовался для
удалённого администрирования различными сетевыми устройствами и программными
серверами. RFC: 5198.

\subsubsection{SIP (Session Initiation Protocol)}
SIP — протокол передачи данных, описывающий способ установления и завершения
пользовательского сеанса связи, включающего обмен мультимедийным содержимым.
RFC: 3261.

\subsubsection{XDR (External Data Representation)}
XDR — международный стандарт передачи данных в Интернете, используемый в различных
RFC для описания типов. XDR позволяет организовать не зависящую от платформы передачу
данных между компьютерами в разнородных сетях. RFC: 4506.

\subsubsection{LPP (Lightweight Presentation Protocol)}
LPP — протокол, используемый для предоставления услуг представления ISO поверх
стеков протоколов на основе TCP/IP. RFC: 1085.

\subsubsection{NDR}
Протокол передачи NDR, также называемый синтаксисом передачи, позволяет вызовам
RPC проходить по сети. Протокол провода определяет проводное представление вызова RPC,
например, порядок отправки элементов данных, выравнивание данных в проводной сети,
дополнительную информацию, включенную в данные, и другие проблемы. RFC: 3464.

\subsubsection{PAP (Password Authentication Protocol)}
PAP — протокол простой проверки подлинности, предусматривающий отправку имени
пользователя и пароля на сервер удалённого доступа открытым текстом. Протокол
аутентификации PAP используется в протоколе PPP для предоставления пользователям
доступа к серверным ресурсам. RFC: 2865.

\subsubsection{L2TP (Layer 2 Tunneling Protocol)}
L2TP — в компьютерных сетях туннельный протокол, использующийся для поддержки
виртуальных частных сетей. Главное достоинство L2TP состоит в том, что этот протокол
позволяет создавать туннель не только в сетях IP, но и в таких, как ATM, X.25 и Frame Relay.
RFC: 2661.

\subsubsection{SMPP (Short Message Peer-to-Peer)}
SMPP — одноранговая передача коротких сообщений. Является открытым протоколом в
телекоммуникационной отрасли, который разработан специально, чтобы обеспечить гибкий
интерфейс для обмена SMS-сообщениями между прикладными SMS-платформами,
маршрутизаторами и центрами службы коротких сообщений. RFC: 6118.

\subsubsection{SDP (Session Description Protocol)}
SDP — сетевой протокол прикладного уровня (уровень приложений в стеке TCP/IP),
предназначенный для описания сессии передачи потоковых данных, включая телефонию,
Интернет-радио, приложения мультимедии. Сессия SDP может реализовывать несколько
потоков данных. RFC: 4566.

\subsubsection{SCP (Secure Copy Protocol)}
SCP — утилита и протокол копирования файлов между компьютерами, использующий, в
отличие от утилиты RCP, в качестве транспорта не RSH, а шифрованный SSH. Сходная по
функционалу утилита — SFTP. RFC: 2995.

\subsubsection{PPTP (Point-to-Point Tunneling Protocol)}
PPTP — туннельный протокол типа точка-точка, позволяющий компьютеру устанавливать
защищённое соединение с сервером за счёт создания специального туннеля в стандартной,
незащищённой сети. PPTP помещает кадры PPP в IP-пакеты для передачи по глобальной IP-сети
(к примеру, Интернет). RFC: 2637.

\endinput
